% FILE: 03_architecture_and_primitives.tex
% Project: adversarial
% Thesis Role: adversarial-testing-and-hostile-assumption-verification

\section{Architecture and Primitives}
\label{sec:adversarial-architecture}

\subsection{Core Objects}
\label{sec:adversarial-architecture-objects}

The adversarial simulation engine (\texttt{simulate\_laundering\_spoofing.py}) is organized
around five core objects.

\textbf{IngestionBoundary} is the zero-trust ingestion boundary class. Its \texttt{validate\_log}
method implements five detection rules in sequence: \texttt{reject\_unhashed\_dataframe},
\texttt{check\_signature\_authenticity}, \texttt{check\_cryptographic\_chain},
\texttt{check\_timestamp\_monotonicity}, and \texttt{check\_impossible\_velocity}. The method
returns a structured verdict dictionary containing \texttt{evaluation\_id},
\texttt{analyzed\_log\_hash\_sha256}, \texttt{verdict} (ACCEPTED or REJECTED),
\texttt{detection\_rules} (audit trail of which rules fired), and on rejection, a
\texttt{refusal\_reason} block identifying the failed rule, trace identifier, and detailed
error message.

\textbf{LaunderingRefusalError} is the custom exception raised when any detection rule fails.
Its fields are \texttt{rule\_failed}, \texttt{trace\_id}, and \texttt{detailed\_error}. The
exception is caught within \texttt{validate\_log} and serialized into the
\texttt{refusal\_reason} block, ensuring that every rejection is fully documented.

\textbf{Wasm4pmBridge} is the ctypes FFI bridge to \texttt{libwasm4pm.dylib}. It wraps five
library functions: \texttt{wasm\_init} (arena initialization with memory ceiling),
\texttt{wasm\_alloc} (heap allocation), \texttt{wasm\_parse\_and\_query} (OCEL log parsing
and OCPQ query evaluation), \texttt{wasm\_shred\_heap} (four-pass zeroization of linear
memory), and \texttt{wasm\_verify\_jcs\_signature} (Ed25519 signature verification over
canonical JSON). Error codes are typed constants: \texttt{ERR\_CONFORMANCE\_VIOLATION =
0xFB03} (signature/conformance failures) and \texttt{ERR\_LIFECYCLE\_VIOLATION = 0xFB05}
(out-of-bounds pointer safety).

\textbf{GENUINE\_SYSTEM\_KEY} and \textbf{FORGED\_SYSTEM\_KEY} form the authority key pair
used in signature verification scenarios. The genuine key simulates an ERP system secret;
the forged key simulates an unauthorized adversary attempting to produce structurally valid
but cryptographically invalid traces.

\subsection{Admissible Transitions}
\label{sec:adversarial-architecture-transitions}

The hash chain primitive defines admissible event transitions. The chain link is computed as:
\[
  H_j = \text{SHA256}(P_j \,\|\, H_{j-1} \,\|\, S_j)
\]
where $P_j$ is the canonical JSON serialization of the event payload (deterministic,
\texttt{sort\_keys=True}), $H_{j-1}$ is the previous chain link, and $S_j$ is the
HMAC-SHA256 signature of $P_j$ under the system key. This construction means that any
modification to a payload---including a single-byte timestamp shift---invalidates the
signature, which in turn invalidates the chain link, which propagates the invalidation to all
subsequent events. Event deletion breaks the chain at the point of deletion, producing a
detectable discontinuity between the expected and actual hash links.

Temporal admissibility requires $T_{\text{ns}}(P_j) \geq T_{\text{ns}}(P_{j-1})$ (monotonicity)
and $\Delta T \geq 10^9$ nanoseconds (impossible-velocity threshold: one second minimum
between events). These bounds are operationally motivated: sub-second event sequences in
enterprise process logs indicate clock manipulation, not genuine process execution.

Witness lattice admissibility is governed by the three structural laws in
\texttt{sources/wasm4pm-compat/adversarial-type-law.md}: strict injectivity (Law I, the
Chatman Constraint), exhaustive typestate closure (Law II), and cryptographic isolation of
witness contexts across epochs via domain separators (Law III: $\text{Hash}_{\text{domain}} =
H(\text{DomainID}_{\text{epoch}} \,\|\, \text{WitnessData})$).

\subsection{Boundaries}
\label{sec:adversarial-architecture-boundaries}

The adversarial architecture defines three explicit boundaries. The \textbf{ingestion
boundary} is the Python \texttt{IngestionBoundary} class, which rejects any input that is a
raw pandas DataFrame (no cryptographic envelope), a malformed or missing envelope, or a trace
with a broken hash chain, invalid signature, non-monotonic timestamps, or impossible velocity.

The \textbf{FFI boundary} is the WebAssembly linear memory interface exposed by
\texttt{Wasm4pmBridge}. It enforces out-of-bounds pointer safety (Scenario 8: offset
\texttt{0xFFFFFFFF} returns \texttt{ERR\_LIFECYCLE\_VIOLATION}) and post-decommissioning
memory zeroization (Scenario 9: four-pass \texttt{wasm\_shred\_heap} zeroes all sentinel
bytes before deallocation).

The \textbf{standards audit boundary} is the \texttt{validate\_standards.py} script, which
enforces that all six standards placement files contain no placeholder terms (TODO, FIXME,
stub, etc.), no broken file links, no malformed JSON blocks, and a complete Trans-Standard
Conversions section with a \texttt{LossReport} schema and \texttt{witness\_signature} field.

\subsection{Failure Modes Refused}
\label{sec:adversarial-architecture-failure-modes}

The architecture explicitly refuses four failure modes. First, \textbf{raw DataFrame
injection}: any \texttt{pd.DataFrame} input is immediately rejected with verdict REJECTED and
rule \texttt{reject\_unhashed\_dataframe}, without further processing. Second,
\textbf{signature forgery}: traces manufactured under a forged key fail
\texttt{check\_signature\_authenticity} via \texttt{hmac.compare\_digest}, which is
constant-time and immune to timing side-channels. Third, \textbf{event deletion}: removing
an event from a hash-chained trace breaks the chain at the deletion point, which
\texttt{check\_cryptographic\_chain} detects. Fourth, \textbf{phantom object injection}:
the OCEL adversarial gravity document (\texttt{standards/ocel-adversarial-gravity.md},
Section 2.2) identifies object-to-event desynchronization as a compliant-processor
rejection criterion---dynamic attribute updates not synchronized with active event-to-object
links must be refused.
